% 中文书籍模板(ctexbook)
% 编译方式: 使用 XeLaTeX 或 LuaLaTeX 编译(推荐 XeLaTeX)。
% 该模板默认自动选用系统可用的中文字体,如需指定可加选项,如:
%   \documentclass[11pt,a4paper,fontset=fandol]{ctexbook}
\documentclass[11pt,a4paper]{ctexbook}
\usepackage[a4paper,margin=1.1in]{geometry}
\usepackage{graphicx}
\usepackage{lipsum}
\usepackage[hidelinks]{hyperref}

\title{ROS2操作系统学习记录}
\author{魏舟}
\date{\today}

\begin{document}

\frontmatter
\maketitle

\tableofcontents

\chapter*{前言}
\addcontentsline{toc}{chapter}{前言}


\mainmatter

\part{基础篇}

\chapter{Unbuntu常用指令}
\section{打开终端}
第一件事是学会打开终端。常用的终端打开方式无外乎两种：
\begin{itemize}
    \item 使用快捷键 Ctrl+Alt+T
    \item 文件夹中右击打开
\end{itemize}
\section{另一个小节}
继续书写正文占位内容。章节标题、目录、页眉页脚都会自动使用中文排版规则。

\chapter{故事继续}
这里是第二个章节的内容占位,展示章与章之间的衔接方式。
新起一段继续书写。\lipsum[4-5]

\part{进阶篇}

\chapter{进入高潮}
这一部分用于演示后续章节的组织方式,例如按主题或难度划分内容。

\backmatter
\chapter{后记}
结语与补充说明可以放在后置部分,例如致谢、术语表或附录说明。

\end{document}
